% Source-faithful assembly: Mumford, Abelian Varieties, Chapter II,
% Sections 4--6 (pages 50--66 of the source PDF). Running headers and
% pedagogical questions/examples are omitted.

\section{Definition of abelian varieties}
\label{mumford-source-II4}
\dcref{M:ch2:II.4}

\begin{definition}[Abelian variety]
\label{mumford-def-abelian-variety}
\dcref{M:ch2:II.4}
An abelian variety $X$ is a complete algebraic variety over an algebraically
closed field $k$ with a group law $m:X\times X\to X$ such that $m$ and the
inverse map are morphisms of varieties.
\source{mumford-abelian-varieties:page-0050}\source{mumford-abelian-varieties:page-0051}
\end{definition}

\begin{proposition}[Basic structure]
\label{mumford-source-basic-structure}
\dcref{M:ch2:II.4}
If $g=\dim X$, then
\[
 X_n\simeq(\mathbf Z/n\mathbf Z)^{2g}\quad(\operatorname{char}k\nmid n),
 \qquad X_{p^m}\simeq(\mathbf Z/p^m\mathbf Z)^i
\]
for $p=\operatorname{char}k$, where $0\le i\le g$; $X$ is commutative and
divisible. Moreover
\[
 H^q(X,\Omega^p)\simeq\bigwedge^p H^1(X,\Omega^1)\otimes_k
 \bigwedge^q H^1(X,\mathcal O_X),
\quad \dim H^1(X,\mathcal O_X)=\dim H^0(X,\Omega^1)=g.
\]
\source{mumford-abelian-varieties:page-0050}
\end{proposition}

\begin{lemma}[Translations]
\label{mumford-lem-abelian-translations}
\dcref{M:ch2:II.4}
For every $a\in X$, the translation $T_a(x)=a+x$ is an automorphism of $X$;
its inverse is $T_{-a}$.  Consequently, local geometric properties are
preserved by translation.
\source{mumford-abelian-varieties:page-0052}\source{mumford-abelian-varieties:page-0053}
\uses{mumford-def-abelian-variety}
\end{lemma}
\begin{proof}
The group identities give $T_{-a}\circ T_a=T_a\circ T_{-a}=\operatorname{id}$,
and both maps are morphisms because the group law and inverse are morphisms.
\end{proof}

\begin{proposition}[Fundamental group and free finite quotients]
\label{mumford-source-II4-roadmap}
\dcref{M:ch2:II.4}
If $g=\dim X$ and $p=\operatorname{char}k$, then the algebraic fundamental group
has the expected form
\[
 \pi_1(X)\simeq\prod_\ell(\mathbf Z_\ell)^{2g}\quad(p=0),\qquad
 \pi_1(X)\simeq\prod_{\ell\ne p}(\mathbf Z_\ell)^{2g}\times\mathbf Z_p^i
 \quad(p>0),
\]
where $i$ is the $p$-rank.  More generally, if a finite group acts freely on a
variety $Y$ with quotient $X$, then some $n>0$ and homomorphisms fit into
\[
\begin{array}{ccccc}&&Y&&\\[-2pt]&\swarrow&&\searrow&\\[-2pt]
X&&\xrightarrow{\,n_X\,}&&X.
\end{array}
\]
The variety $Y$ carries an abelian-variety structure for which both maps are
homomorphisms.
\source{mumford-abelian-varieties:page-0050}\source{mumford-abelian-varieties:page-0051}
% Source states this roadmap result on pp. 50--51 but gives no proof there.
\notready
\end{proposition}

\begin{lemma}[Everywhere nonsingular]
\label{mumford-source-nonsingular}
\dcref{M:ch2:II.4}
An abelian variety is everywhere nonsingular.
\source{mumford-abelian-varieties:page-0052}
\end{lemma}
\begin{proof}
There is a nonsingular point $x_0$. For $x\in X$, the translation morphism
$T_{(x,x_0^{-1})}(y)=x\cdot x_0^{-1}\cdot y$ is an automorphism carrying
$x_0$ to $x$, so $x$ is nonsingular.
\end{proof}

\begin{theorem}[First commutativity proof]
\label{mumford-source-commutativity-I}
\dcref{M:ch2:II.4}
An abelian variety is commutative.
\source{mumford-abelian-varieties:page-0052}
\end{theorem}
\begin{proof}
For $x\in X$, conjugation $C_x(y)=xyx^{-1}$ induces an automorphism
$C^*_{x,n}$ of $\mathcal O_{X,e}/\mathfrak m_{X,e}^n$. The map
$X\to\operatorname{Aut}(\mathcal O_{X,e}/\mathfrak m_{X,e}^n)$ is a morphism
to an affine variety, hence constant because $X$ is complete and connected.
It equals the identity at $e$. Thus $C^*_{x,n}=1$ for every $n$; since
$\bigcap_n\mathfrak m_{X,e}^n=0$, $C_x$ is the identity near $e$, hence on
the irreducible variety $X$.
\end{proof}

We write the group law additively. Let $T_x(y)=x+y$ and let $n_X$ denote
multiplication by $n$.
\source{mumford-abelian-varieties:page-0052}\source{mumford-abelian-varieties:page-0053}

\begin{proposition}[Invariant differentials]
\label{mumford-source-invariant-differentials}
\dcref{M:ch2:II.4}
If $T=T_{X,0}$ and $\Omega_0=T^*_{X,0}$, then
\[
 \Omega_0\otimes_k\mathcal O_X\xrightarrow{\sim}\Omega_X^1.
\]
Every regular differential form on $X$ is translation invariant.
\source{mumford-abelian-varieties:page-0053}
\end{proposition}
\begin{proof}
For $\theta\in\Omega_0$, define $\omega_\theta$ by
$(\omega_\theta)_x=T^*_{-x}(\theta)$. This is regular and gives the map.
At each $x$, $T^*_{-x}$ identifies the fiber with $\Omega_0$, so Nakayama's
lemma proves it is an isomorphism. Since $H^0(X,\mathcal O_X)=k$, the claim
about global forms follows.
\end{proof}

\begin{theorem}[Prime-to-characteristic multiplication]
\label{mumford-prop-multiplication-n}
\dcref{M:ch2:II.4}
For every $n$ not divisible by $\operatorname{char}k$, $n_X$ is surjective.
\source{mumford-abelian-varieties:page-0053}\source{mumford-abelian-varieties:page-0054}
\end{theorem}
\begin{proof}
The tangent space at $(0,0)$ splits as $T\oplus T$, and the differential of
$m$ is addition: $d(m)(t_1,t_2)=t_1+t_2$, since the restrictions to the two
summands are the identity. Induction gives $(dn_X)_0=n$ on $T$. Thus it is an
isomorphism when $\operatorname{char}k\nmid n$. If $n_X$ were not surjective,
the dimension theorem would give a nonzero tangent vector to
$n_X^{-1}(0)$ at $0$, contradicting injectivity of $(dn_X)_0$.
\end{proof}

\begin{lemma}[Rigidity Lemma, Form I]
\label{mumford-lem-rigidity-form-I}
\dcref{M:ch2:II.4}
Let $X$ be complete, $Y,Z$ varieties, and $f:X\times Y\to Z$ a morphism such
that $f(X\times\{y_0\})=\{z_0\}$. Then $f=g\circ p_2$ for a morphism
$g:Y\to Z$.
\source{mumford-abelian-varieties:page-0054}
\end{lemma}
\begin{proof}
Choose $x_0\in X$ and set $g(y)=f(x_0,y)$. Let $U$ be an affine neighborhood
of $z_0$, $F=Z-U$, and $G=p_2(f^{-1}(F))$. Since $X$ is complete, $G$ is
closed; since $y_0\notin G$, $V=Y-G$ is nonempty open. For $y\in V$, the
complete variety $X\times\{y\}$ maps into affine $U$, hence to one point.
Thus $f(x,y)=f(x_0,y)=g(p_2(x,y))$ on a nonempty open, and therefore everywhere.
\end{proof}

\begin{corollary}[Morphisms are homomorphisms up to translation]
\label{mumford-cor-morphism-translation}
\dcref{M:ch2:II.4}
If $X,Y$ are abelian varieties and $f:X\to Y$ is a morphism, then
$f(x)=h(x)+a$ for a homomorphism $h$ and $a\in Y$.
\source{mumford-abelian-varieties:page-0054}
\uses{mumford-lem-rigidity-form-I}
\end{corollary}
\begin{proof}
Replace $f$ by $f-f(0)$. The morphism
$\phi(x,y)=f(x+y)-f(y)-f(x)$ vanishes on both coordinate axes. Rigidity gives
$\phi=0$, so $f$ is a homomorphism.
\end{proof}

\begin{corollary}[Second commutativity proof]
\label{mumford-cor-abelian-commutative}
\dcref{M:ch2:II.4}
An abelian variety is commutative.
\source{mumford-abelian-varieties:page-0055}
\uses{mumford-cor-morphism-translation}
\end{corollary}
\begin{proof}
The inverse map is a homomorphism by the preceding corollary. Hence
$(xy)^{-1}=x^{-1}y^{-1}=y^{-1}x^{-1}$, so $X$ is commutative.
\end{proof}

\begin{corollary}[Linearity of maps into an abelian variety]
\label{mumford-cor-pointed-hom-linear}
\dcref{M:ch2:II.4}
For pointed complete varieties $S,T$ and an abelian variety $X$,
\[
 \Hom(S,X)\times\Hom(T,X)\xrightarrow{\sim}\Hom(S\times T,X),
 \qquad(f,g)\mapsto[(s,t)\mapsto f(s)+g(t)].
\]
\source{mumford-abelian-varieties:page-0055}
\end{corollary}
\begin{proof}
Injectivity follows by setting either variable equal to its base point. Given
$h$, set $f(s)=h(s,t_0)$, $g(t)=h(s_0,t)$ and
$k(s,t)=h(s,t)-f(s)-g(t)$. Rigidity applied to $k$ gives $k=0$.
\end{proof}

\begin{theorem}[Recognition criterion]
\label{mumford-thm-unital-composition}
\dcref{M:ch2:II.4}
Let $X$ be complete, $e\in X$, and $m:X\times X\to X$ satisfy
$m(x,e)=m(e,x)=x$. Then $m$ is an abelian-variety group law with identity $e$.
\source{mumford-abelian-varieties:page-0055}\source{mumford-abelian-varieties:page-0056}
\end{theorem}
\begin{proof}
Write $xy=m(x,y)$ and define $\psi(x,y)=(xy,y)$. Since
$\psi^{-1}(e,e)=\{(e,e)\}$, the dimension theorem and completeness make
$\psi$ surjective. Let $\Gamma'$ be $\{(x,y):xy=e\}$, choose an irreducible
component $\Gamma$ projecting onto the second factor, and then onto the first
factor by the dimension theorem. Rigidity applied to
$\phi((x',x),y)=x'(xy)$ gives
\[
 x'(xy)=y\quad((x',x)\in\Gamma).
 \tag{1}
\]
Choosing $(x'',x')\in\Gamma$ and using (1) shows also $xx'=e$.

Apply rigidity to $\chi(((x',x),y,z))=x'((xy)z)$ to obtain
$x'((xy)z)=yz$. Multiplying on the left by $x'$ and using (1) gives
$(x'y)z=x'(yz)$. Since $x'$ is arbitrary, multiplication is associative.
Translations are then automorphisms, so $X$ is nonsingular; $\psi$ is
bijective. Its tangent map at $(e,e)$ is an isomorphism, so it is inseparable
and, by Zariski's Main Theorem, an isomorphism. Its inverse is
$(x,y)\mapsto(xy^{-1},y)$, proving that inversion is a morphism.
\end{proof}

\section{Cohomology and base change}
\label{mumford-source-II5}
\dcref{M:ch2:II.5}

\begin{theorem}[Coherence of proper direct images]
\label{mumford-thm-proper-direct-image-coherent}
\dcref{M:ch2:II.5}
If $f:X\to Y$ is proper between locally noetherian preschemes and $\mathcal F$
is coherent, then $R^pf_*\mathcal F$ is coherent for all $p\ge0$.
\source{mumford-abelian-varieties:page-0057}
% Source explicitly takes this coherence theorem as a basic result; no proof is given.
\notready
\end{theorem}

\begin{definition}[Flat over the base]
\label{mumford-def-f-flat}
\dcref{M:ch2:II.5}
A quasicoherent $\mathcal F$ on $X$ is $f$-flat if every $\mathcal F_x$ is
$\mathcal O_{Y,f(x)}$-flat; equivalently, sections over affine $U$ mapping into
affine $V$ are flat $\mathcal O_Y(V)$-modules.
\source{mumford-abelian-varieties:page-0057}
\end{definition}

\begin{theorem}[Cohomology and base change complex]
\label{mumford-thm-grothendieck-complex}
\dcref{M:ch2:II.5}
For proper $f:X\to Y$ with $Y=\Spec A$, $X$ noetherian, and coherent
$\mathcal F$ flat over $Y$, there is a finite complex $K^\bullet$ of finitely generated
projective $A$-modules such that, functorially for every $A$-algebra $B$,
\[
 H^p(X\times_Y\Spec B,\mathcal F\otimes_A B)
 \simeq H^p(K^\bullet\otimes_A B).
\]
\source{mumford-abelian-varieties:page-0057}\source{mumford-abelian-varieties:page-0058}\source{mumford-abelian-varieties:page-0059}\source{mumford-abelian-varieties:page-0060}
\end{theorem}
\begin{proof}
Choose a finite affine covering and its Cech complex $C^\bullet$. It is a
finite complex of flat $A$-modules and base change identifies its cohomology
with the fiber cohomology. Lemma~\ref{mumford-source-complex-model} supplies
a finite complex $K^\bullet$ of finite modules, free in positive degrees,
quasi-isomorphic to $C^\bullet$, with $K^0$ flat. Since $A$ is noetherian,
$K^0$ is projective.

For completeness, the construction descends on $m$: choose a free module
mapping onto the kernel of $\ker\partial^{m+1}\to H^{m+1}(C^\bullet)$ and a
free module mapping onto $H^m(C^\bullet)$, lift the latter to cycles, and
take their direct sum. At $m=-1$ quotient $K^0$ by
$\ker\partial^0\cap\ker\phi_0$. The mapping cylinder
$L^p=K^p\oplus C^{p-1}$ is exact; if all $C^p$ are flat, the short exact
sequences $0\to Z^p\to L^p\to Z^{p+1}\to0$ show $K^0=L^0$ is flat.

Lemma~\ref{mumford-lem-flat-quasi-iso-base-change} applied to $K^\bullet\to C^\bullet$
shows that base change preserves the quasi-isomorphism, proving the claim.
\end{proof}

\begin{lemma}[Finite complex model]
\label{mumford-source-complex-model}
\dcref{M:ch2:II.5}
Let $C^\bullet$ be bounded over a noetherian ring $A$, with finitely generated
cohomology. There is a bounded complex $K^\bullet$ of finite modules, free in
positive degrees, mapping quasi-isomorphically to $C^\bullet$; if $C^p$ are
flat, $K^0$ is flat.
\source{mumford-abelian-varieties:page-0058}\source{mumford-abelian-varieties:page-0059}\source{mumford-abelian-varieties:page-0060}
\end{lemma}
\begin{proof}
Construct $K^\bullet$ by descending induction. Put $K^p=0$ for $p>n$.
Assume the maps have been constructed above degree $m$ so that they induce
cohomology isomorphisms above degree $m+1$ and a surjection
$\ker\partial^{m+1}\to H^{m+1}(C^\bullet)$. For $m\ge0$, let
$B^{m+1}$ be the kernel of this surjection. Choose a finite free ${K'}^m$
surjecting onto $B^{m+1}$, and a finite free ${K''}^m$ surjecting onto
$H^m(C^\bullet)$; lift the latter map to cycles. Set
$K^m={K'}^m\oplus{K''}^m$, let the differential vanish on ${K''}^m$ and
map ${K'}^m$ onto $B^{m+1}$, and lift this map to $C^m$. The resulting map
of complexes has the required cohomology properties. For $m=-1$, replace
$K^0$ by $K^0/(\ker\partial^0\cap\ker\phi_0)$.

If all $C^p$ are flat, form the mapping cylinder
$L^p=K^p\oplus C^{p-1}$ with
$\partial(x,0)=(\partial x,\phi(x))$ and
$\partial(0,y)=(0,-\partial y)$. The exact sequence of complexes with the
shifted $C^\bullet$ shows $H^p(L^\bullet)=0$. Hence
$0\to K^0=L^0\to L^1\to\cdots\to L^{n+1}\to0$ is exact; its other terms
are flat, so $K^0$ is flat.
\end{proof}

\begin{lemma}[Flat complexes commute with arbitrary base change]
\label{mumford-lem-flat-quasi-iso-base-change}
\dcref{M:ch2:II.5}
If finite complexes of flat $A$-modules are quasi-isomorphic, then their
cohomologies remain isomorphic after tensoring with every $A$-algebra $B$.
\source{mumford-abelian-varieties:page-0060}
\end{lemma}
\begin{proof}
Construct the mapping cylinder $L^\bullet$ of the quasi-isomorphism. It is an
exact finite complex of flat modules. Its cycles $Z^p$ are flat, and
\[
0\to Z^p\to L^p\to Z^{p+1}\to0
\]
is exact. Tensoring with $B$ remains exact, so $L^\bullet\otimes_A B$ is
exact. It is the mapping cylinder after base change; the cohomology sequence
therefore gives the required isomorphisms.
\end{proof}

\begin{corollary}[Semicontinuity and Euler characteristic]
\label{mumford-thm-semicontinuity}
\dcref{M:ch2:II.5}
For proper $f:X\to Y$ and coherent $\mathcal F$ flat over $Y$, the function
$y\mapsto\dim_{k(y)}H^p(X_y,\mathcal F_y)$ is upper semicontinuous, while
$y\mapsto\chi(\mathcal F_y)$ is locally constant.
\source{mumford-abelian-varieties:page-0061}
\end{corollary}
\begin{proof}
Locally choose the free complex $K^\bullet$ and write $d^p:K^p\to K^{p+1}$.
Then
\[
\dim H^p=\dim K^p\otimes k(y)-\dim\operatorname{Im}(d^p\otimes k(y))
 -\dim\operatorname{Im}(d^{p-1}\otimes k(y)).
\]
The alternating sum is locally constant. The image rank is lower
semicontinuous because its rank-drop locus is cut out by the entries of the
induced exterior-power matrix, proving upper semicontinuity.
\end{proof}

\begin{corollary}[Base-change criterion]
\label{mumford-thm-base-change-criterion}
\dcref{M:ch2:II.5}
If $Y$ is reduced and connected, constancy of
$y\mapsto\dim H^p(X_y,\mathcal F_y)$ is equivalent to $R^pf_*\mathcal F$
being locally free and having fibers $H^p(X_y,\mathcal F_y)$. Under these
conditions the same base-change map is an isomorphism in degree $p-1$.
\source{mumford-abelian-varieties:page-0061}\source{mumford-abelian-varieties:page-0062}\source{mumford-abelian-varieties:page-0063}
\end{corollary}
\begin{proof}
First, on a reduced scheme, a coherent sheaf with constant fiber dimension
$r$ is locally free of rank $r$: lift generators at a point, apply Nakayama,
and observe that the kernel lies in $\mathfrak m_y\mathcal O_y^r$ at every
point, hence vanishes on a reduced scheme. Second, if
$\phi:\mathscr F\to\mathfrak D$ is a map of locally free sheaves with
locally constant image rank, the preceding result makes the cokernel locally
free. Thus the image is projective and both source and target split, with
matrix $\left[\begin{smallmatrix}0&\text{isom.}\\0&0\end{smallmatrix}\right]$.
Applying the splitting to $d_{p-1}$ and $d_p$ in $K^\bullet$ gives projective
modules $B_p,H_p,K'_p$ with
$B_p\oplus H_p=\ker d_p$ and $d_p:K'_p\simeq B_{p+1}$. Hence
$H^p(K^\bullet\otimes B)=H_p\otimes B=H^p(K^\bullet)\otimes B$ and likewise
in degree $p-1$.
\end{proof}

\begin{corollary}[Vanishing in the preceding degree]
\label{mumford-source-base-change-vanishing}
\dcref{M:ch2:II.5}
If $H^p(X_y,\mathcal F_y)=0$ for all $y$, then
$R^{p-1}f_*\mathcal F\otimes k(y)\simeq H^{p-1}(X_y,\mathcal F_y)$ for all $y$.
\source{mumford-abelian-varieties:page-0064}
\end{corollary}
\begin{proof}
Locally write the relevant part of $K^\bullet$ as
$K^{p-1}\to K^p\to K^{p+1}$. Exactness after tensoring with $k(y)$ permits
a splitting $K^p\otimes k(y)=\overline W_1\oplus\overline W_2$, where the
first summand is the image and the second maps injectively. After localizing
at a function nonzero at $y$, lift bases to a splitting $K^p=W_1\oplus W_2$
with $W_1\subset\operatorname{Im}d^{p-1}$; injectivity on $W_2$ persists,
so $W_1=\operatorname{Im}d^{p-1}$. The surjection onto $W_1$ splits, giving
$K^{p-1}=\ker d^{p-1}\oplus W_1$ and exact sequences
\[
 K^{p-2}\to\ker d^{p-1}\to H^{p-1}(X,\mathcal F)\to0,
\quad
 K^{p-2}\otimes k(y)\to\ker d^{p-1}\otimes k(y)\to
 H^{p-1}(X_y,\mathcal F_y)\to0.
\]
This proves the asserted base-change isomorphism.
\end{proof}

\begin{corollary}[Higher vanishing]
\label{mumford-source-higher-vanishing}
\dcref{M:ch2:II.5}
If $R^kf_*\mathcal F=0$ for $k>k_0$, then $H^k(X_y,\mathcal F_y)=0$ for
$k>k_0$ and all $y$.
\source{mumford-abelian-varieties:page-0064}
\end{corollary}
\begin{proof}
Apply the preceding vanishing-in-the-preceding-degree corollary and use
decreasing induction on $k_0$.
\end{proof}

\begin{corollary}[Flat base change]
\label{mumford-cor-flat-base-change}
\dcref{M:ch2:II.5}
If $B$ is flat over $A$, then
\[
 H^p(X\times_Y\Spec B,\mathcal F\otimes_A B)\simeq H^p(X,\mathcal F)\otimes_A B.
\]
\source{mumford-abelian-varieties:page-0064}\source{mumford-abelian-varieties:page-0065}
\end{corollary}
\begin{proof}
For a flat algebra, cohomology of any complex commutes with tensor product.
\end{proof}

\begin{theorem}[Seesaw theorem, provisional form]
\label{mumford-thm-seesaw}
\dcref{M:ch2:II.5}
For complete $X$, any variety $T$, and a line bundle $L$ on $X\times T$, the
set of $t$ for which $L|_{X\times\{t\}}$ is trivial is closed. On this closed
set $T_1$, $L|_{X\times T_1}\simeq p_2^*M$ for a line bundle $M$ on $T_1$.
\source{mumford-abelian-varieties:page-0065}\source{mumford-abelian-varieties:page-0066}
\uses{mumford-thm-base-change-criterion}
\end{theorem}
\begin{proof}
A line bundle $M$ on complete $X$ is trivial iff both $H^0(X,M)$ and
$H^0(X,M^{-1})$ are nonzero: dual nonzero sections compose to a nonzero
constant, giving inverse isomorphisms. Thus $T_1$ is the intersection of the
two closed loci supplied by semicontinuity. After replacing $T$ by $T_1$,
all fibers are trivial and have one-dimensional $H^0$. Base-change criterion
gives $p_{2*}\underline L=\underline M$ invertible and the natural map
$p_2^*\underline M\to\underline L$ is an isomorphism.
\end{proof}

\section{The theorem of the cube: I}
\label{mumford-source-II6}
\dcref{M:ch2:II.6}

\begin{theorem}[Theorem of the cube]
\label{mumford-thm-cube-varieties}
\dcref{M:ch2:II.6}
Let $X,Y$ be complete varieties, $Z$ any variety, with base points
$x_0,y_0,z_0$. If a line bundle $L$ on $X\times Y\times Z$ restricts
trivially to each of $\{x_0\}\times Y\times Z$,
$X\times\{y_0\}\times Z$, and $X\times Y\times\{z_0\}$, then $L$ is trivial.
\source{mumford-abelian-varieties:page-0066}
% The source proof is assembled in ch02-source-part3.tex from pp. 67--70.
\end{theorem}

\begin{remark}[Functorial order]
\label{mumford-source-functor-order}
\dcref{M:ch2:II.6}
For a contravariant functor $T$ on complete varieties, define projections
$\pi_i$ and base-point inclusions $\sigma_i$ by omitting and restoring the
$i$th factor. Then
\[
 \alpha_T^n(\xi_0,\ldots,\xi_n)=\sum_{i=0}^n\pi_i^*(\xi_i),
 \qquad
 \beta_T^n(\eta)=(\sigma_0^*\eta,\sigma_2^*\eta,\ldots,\sigma_n^*\eta),
\]
and induction gives $T(\prod X_i)=\operatorname{Im}\alpha\oplus\ker\beta$.
The functor has order $n$ when $\alpha$ is surjective (equivalently $\beta$
injective). The cube theorem, when $Z$ is complete, says that $\Pic X$ is
quadratic.
\source{mumford-abelian-varieties:page-0066}
\end{remark}
